\section{Introducción: por qué se necesita el modelado de secuencias}

Esta clase es la segunda de MIT 6.S191 (Introduction to Deep Learning), impartida por Ava Amini. En la clase anterior, Alexander Amini presentó los fundamentos de las redes neuronales, incluyendo el perceptrón (Perceptron), la red de propagación hacia adelante (Feedforward Network) y el algoritmo de retropropagación (Backpropagation). A partir de estos fundamentos, esta clase explica de manera sistemática cómo aplicar las redes neuronales al procesamiento de \textbf{datos secuenciales} (Sequential Data), y abarca tres temas centrales:

\begin{enumerate}
    \item \textbf{Recurrent Neural Networks (RNNs)}: los principios básicos y el mecanismo interno de las redes neuronales recurrentes
    \item \textbf{Los retos de entrenar una RNN}: el problema del desvanecimiento/explosión del gradiente y arquitecturas mejoradas como la LSTM
    \item \textbf{El mecanismo de Attention y el Transformer}: el principio del mecanismo de autoatención y su aplicación en los modelos de lenguaje grandes modernos
\end{enumerate}

\begin{importantbox}{La idea central del modelado de secuencias}
El objetivo central del modelado de secuencias es: dada la información histórica de una secuencia, predecir su estado futuro. Esta es la base de numerosas aplicaciones de IA, como los modelos de lenguaje grandes (LLM), el reconocimiento de voz y la traducción automática.
\end{importantbox}

\subsection{Introducción intuitiva: predecir el movimiento de una pelota}

La ponente introduce el concepto de modelado de secuencias con un ejemplo ilustrativo: supongamos que hay una pelota en un espacio bidimensional, y nuestra tarea es predecir la dirección de su próximo movimiento.

\begin{figure}[H]
\centering
\includegraphics[width=0.7\textwidth]{slides-images/slide-03.jpg}
\caption{Solo con la posición actual de la pelota no se puede determinar la dirección de su movimiento\protect\footnotemark}
\end{figure}
\footnotetext{Diapositivas, página 3.}

Si solo se da la posición actual de la pelota, sin ninguna información histórica, nuestra predicción solo puede ser aleatoria: la pelota podría moverse en cualquier dirección. Pero si se da la secuencia de posiciones de la pelota en los pasos de tiempo anteriores, el problema se vuelve mucho más sencillo:

\begin{figure}[H]
\centering
\includegraphics[width=0.7\textwidth]{slides-images/slide-05.jpg}
\caption{Usar la secuencia de posiciones históricas permite predecir con precisión la dirección del movimiento\protect\footnotemark}
\end{figure}
\footnotetext{Diapositivas, página 5.}

Al observar la trayectoria histórica de la pelota de izquierda a derecha, podemos inferir razonablemente que seguirá moviéndose hacia la derecha. Este ejemplo revela la idea central del modelado de secuencias: \textbf{usar la información pasada para predecir el futuro}.

\subsection{Los datos secuenciales están en todas partes}

Los datos secuenciales están presentes en todas partes del mundo real:

\begin{figure}[H]
\centering
\includegraphics[width=0.85\textwidth]{slides-images/slide-07.jpg}
\caption{Tipos de datos secuenciales en el mundo real\protect\footnotemark}
\end{figure}
\footnotetext{Diapositivas, página 7.}

\begin{itemize}
    \item \textbf{Señales de audio}: la forma de onda del habla puede dividirse en una serie de segmentos de onda sonora para procesarse como secuencia
    \item \textbf{Lenguaje natural}: el texto puede dividirse en una secuencia de palabras o caracteres
    \item \textbf{Señales médicas}: por ejemplo, la serie de tiempo de un electrocardiograma (ECG)
    \item \textbf{Datos financieros}: la serie de tiempo del precio de las acciones
    \item \textbf{Secuencias biológicas}: secuencias de ácidos nucleicos de ADN/ARN, secuencias de aminoácidos de proteínas
    \item \textbf{Datos meteorológicos, secuencias de fotogramas de video}, entre otros
\end{itemize}

\subsection{Tipos de tareas del modelado de secuencias}

\begin{knowledgebox}{Los tres patrones básicos del modelado de secuencias}
\begin{itemize}
    \item \textbf{Many-to-One}: entrada de secuencia $\rightarrow$ salida única. Por ejemplo: la clasificación de sentimiento (se ingresa un fragmento de texto y se produce positivo/negativo)
    \item \textbf{One-to-Many}: entrada única $\rightarrow$ salida de secuencia. Por ejemplo: la generación de descripciones de imágenes (se ingresa una imagen y se produce el texto descriptivo)
    \item \textbf{Many-to-Many}: entrada de secuencia $\rightarrow$ salida de secuencia. Por ejemplo: la traducción automática (secuencia en inglés $\rightarrow$ secuencia en chino)
\end{itemize}
\end{knowledgebox}

\begin{figure}[H]
\centering
\includegraphics[width=0.75\textwidth]{slides-images/slide-08.jpg}
\caption{Los distintos tipos de tareas del modelado de secuencias\protect\footnotemark}
\end{figure}
\footnotetext{Diapositivas, página 8.}

\subsection{Resumen de la sección}

El modelado de secuencias es uno de los paradigmas más importantes del aprendizaje profundo, y su núcleo consiste en aprovechar la dependencia temporal/secuencial de los datos. Desde la simple predicción del movimiento de una pelota hasta los modelos de lenguaje complejos, la idea del modelado de secuencias atraviesa todos los aspectos de la IA moderna.

\newpage
